\usepackage{fontspec}
\usepackage[english, russian]{babel}
\usepackage{graphicx}
\usepackage{listings}
\usepackage[a4paper,
            left=3cm,
            right=3cm,
            top=2cm,
            bottom=3cm]{geometry}
\usepackage{amsmath, amssymb}
\usepackage{pgfplots}
    \usetikzlibrary{patterns}
\usepackage{float}

\setmainfont{CMU Serif}
\setsansfont{CMU Sans Serif}
\setmonofont{CMU Typewriter Text}

\usepackage{hyperref}
\hypersetup{
    colorlinks=true,        % цветные ссылки (false - черные рамки)
    linkcolor=black,        % цвет ссылок на разделы
    urlcolor=blue,          % цвет URL-ссылок
    citecolor=green,        % цвет ссылок на литературу
    pdfborder={0 0 0}       % убрать рамки вокруг ссылок
}

% Настройки для листингов кода
\lstset{
    language=Python,
    basicstyle=\ttfamily\small,
    keywordstyle=\color{blue},
    commentstyle=\color{green!60!black},
    stringstyle=\color{red},
    numbers=left,
    numberstyle=\tiny,
    stepnumber=1,
    numbersep=5pt,
    frame=single,
    tabsize=4,
    breaklines=true,
    showstringspaces=false,
    captionpos=b
}

\usepackage{titlesec}

\setcounter{secnumdepth}{4}
\setcounter{tocdepth}{4}

\newcommand{\abs}[1]{\left|#1\right|}